\documentclass[11pt]{article}

\usepackage{amssymb,amsmath,amsfonts,amsthm}
\usepackage{latexsym}
\usepackage{graphics}
\usepackage{indentfirst}
\usepackage{lmodern}
\usepackage{hyperref}
\usepackage{mathrsfs}
\usepackage{array}
\usepackage{mathtools}
\usepackage{graphicx}
\usepackage{setspace}
\usepackage{enumerate}
\usepackage{tikz}
\usepackage{url,color}
\usepackage{xcolor}
\usepackage{chngcntr}
\usepackage{multicol}
\usepackage{comment}
\usepackage{float}
\usepackage{algpseudocode}
\usepackage{subcaption}
\usepackage[ruled]{algorithm2e}
\usepackage{booktabs}
\usepackage{caption}
\usepackage[letterpaper,margin=2cm,top=2.5cm,bottom=2.5cm]{geometry}

\AtBeginEnvironment{algorithm}{\setstretch{1.2}}

\hypersetup{colorlinks=true,citecolor=blue,linkcolor=blue,urlcolor=blue}

%%%%%%%%%%%%%%%%%%%%%%%%%%%%%%%%%%%%%%%%%%%%%%%%%%%%%%%%%%
%% Theorem environments
%%%%%%%%%%%%%%%%%%%%%%%%%%%%%%%%%%%%%%%%%%%%%%%%%%%%%%%%%%
\newtheorem{theorem}{Theorem}[section]
\newtheorem{lemma}{Lemma}[section]
\newtheorem{proposition}{Proposition}[section]
\newtheorem{corollary}{Corollary}[section]
\newtheorem{claim}{Claim}[section]
\newtheorem{conjecture}{Conjecture}[section]
\newtheorem{hypothesis}{Hypothesis}[section]
\newtheorem{assumption}{Assumption}[section]
\newtheorem{remark}{Remark}[section]
\newtheorem{example}{Example}[section]
\newtheorem{problem}{Problem}[section]
\newtheorem{definition}{Definition}[section]
\newtheorem{question}{Question}[section]
\newtheorem{condition}{Condition}[section]

%%%%%%%%%%%%%%%%%%%%%%%%%%%%%%%%%%%%%%%%%%%%%%%%%%%%%%%%%%
%% Math shorthands
%%%%%%%%%%%%%%%%%%%%%%%%%%%%%%%%%%%%%%%%%%%%%%%%%%%%%%%%%%
\newcommand{\N}{\mathbb{N}}
\newcommand{\Z}{\mathbb{Z}}
\newcommand{\E}{\mathbb{E}}
\newcommand{\Q}{\mathbb{Q}}
\newcommand{\R}{\mathbb{R}}
\newcommand{\C}{\mathbb{C}}
\newcommand{\norm}[1]{\left\lVert#1\right\rVert}

%%%%%%%%%%%%%%%%%%%%%%%%%%%%%%%%%%%%%%%%%%%%%%%%%%%%%%%%%%
%% Draft-note macros (REMOVE before submission)
%%%%%%%%%%%%%%%%%%%%%%%%%%%%%%%%%%%%%%%%%%%%%%%%%%%%%%%%%%
\newcommand{\todo}[1]{\textcolor{red}{\textbf{[TODO:}\,#1\textbf{]}}}
\newcommand{\needdata}[1]{\textcolor{orange}{\textbf{[NEED DATA:}\,#1\textbf{]}}}
\newcommand{\draft}[1]{\textcolor{blue}{#1}}
\newcommand{\note}[1]{\textcolor{purple}{\textbf{[NOTE:}\,#1\textbf{]}}}

\allowdisplaybreaks

%%%%%%%%%%%%%%%%%%%%%%%%%%%%%%%%%%%%%%%%%%%%%%%%%%%%%%%%%%
%% Title, authors, abstract
%%%%%%%%%%%%%%%%%%%%%%%%%%%%%%%%%%%%%%%%%%%%%%%%%%%%%%%%%%

\title{A Kernel-Flexible Diffusion Model for Political Redistricting:\\
       Continuous and Integer Local Refinement of ReCom}

\author{Alaittin K\i rt\i \c{s}o\u{g}lu \\
        \small Department of Applied Mathematics, Illinois Institute of Technology \\
        \small Chicago, IL 60616 \\
        \small \texttt{kirtisoglu@hawk.iit.edu}}

\date{\today}

\begin{document}

\maketitle

\begin{abstract}
The Recombination (ReCom) algorithm and its variants dominate Markov chain
Monte Carlo (MCMC) sampling for political redistricting, but every
edge in their spanning-tree split is treated identically and the
proposal mechanism is blind to geographic structure: shared boundary
length, population density, demographic similarity, and administrative
boundaries all carry signal that the chain ignores. We propose a
\emph{kernel-aware local refinement} for redistricting MCMC, formulated
as a single Laplacian energy
\(
  \min_{x}\; \alpha\, x^\top L\, x + \beta\, \|x - x_0\|^2
\)
subject to population balance, where $L$ is a kernel-weighted graph
Laplacian on the merged-pair subgraph and $x_0$ encodes the current
district assignment. The framework supports two solvers for the same
energy: a continuous quadratic-program (QP) relaxation with
Cheeger-sweep rounding and an integer QP (IQP) solved directly as a
mixed-integer quadratic program. Both mix with ReCom for ergodicity.
The framework strictly generalizes SpecReCom (Davies et al., 2025) as
the $(\beta = 0, \text{uniform kernel})$ limit, and subsumes the
discrete cut-min integer program as the integer specialisation.
\draft{Empirical results on Iowa precincts (k=4, $\varepsilon = 0.05$)
show a 23.9\% reduction in cut edges versus Tempered ReCom
($p < 10^{-197}$), with a quantified compactness-coverage trade-off on
a fully enumerated $5 \times 5$ ground truth. A structural diagnostic
of the partition fiber reveals an island-and-corridor macro-structure
whose narrow-throat count is approximately invariant across grid sizes
and balance tolerances; ReCom jumps between islands while spectral local
kernels explore within them, directly motivating the hybrid design.} We
also interpret the QP-diffusion as an implicit Euler step of the graph
heat equation and exhibit its closed-form spectral filter.
\end{abstract}

\tableofcontents
\newpage

%════════════════════════════════════════════════════════════════════════════════
\section{Introduction}\label{sec:intro}
%════════════════════════════════════════════════════════════════════════════════

\subsection{The redistricting problem and ReCom's role}\label{sec:intro-problem}

Political redistricting seeks a partition of a geographic region into
electoral districts satisfying legal, demographic, and geometric
requirements: contiguity, population balance, compactness, and
preservation of communities of interest. Computational approaches model
the problem on the dual adjacency graph $G = (V,E)$ of geographic units
(precincts, census blocks, counties). A valid $k$-district plan is a
partition
\[
  P = \{V_1, \ldots, V_k\}, \qquad \bigsqcup_{i=1}^k V_i = V,
\]
where each $G[V_i]$ is connected and the populations $p(V_i)$ are
approximately equal. Let $\mathcal{P}_k(G)$ denote the space of valid
$k$-district plans. The cardinality of $\mathcal{P}_k(G)$ is enormous:
even four congressional districts over the 99 counties of Iowa admit
roughly $10^{24}$ plans~\cite{fifield2020essential}.

Markov chain Monte Carlo (MCMC) has become the standard tool for
sampling $\mathcal{P}_k(G)$. Early flip chains modify a single
precinct's assignment per step~\cite{fifield2020automated} and mix
slowly. The Recombination (ReCom) algorithm of DeFord, Duchin, and
Solomon~\cite{deford2021recombination} performs multi-precinct updates
by merging two adjacent districts, sampling a uniform random spanning
tree of the merged subgraph, and cutting an edge to recover a balanced
bipartition. ReCom dramatically improved the mixing efficiency of
redistricting MCMC and is now the workhorse of the field, with
reversible~\cite{cannon2022spanning} and multi-scale~\cite{autry2020multiscale}
extensions.

\textbf{ReCom's blind spot.} The spanning-tree-based split is
\emph{kernel-blind}: every edge in the merged subgraph carries the same
implicit weight, and the resulting bipartition is sampled from a
distribution determined by the topology of the merged subgraph alone.
Real geography is not topology-only: shared boundary length, population
density, demographic similarity, and administrative boundaries
all carry signal that should shape the proposed split. The recent
SpecReCom of Davies et al.~\cite{davies2025spectral} partially
addresses this by replacing the random spanning tree with a Fiedler-vector
spectral cut, but it does not admit kernel weighting at the
proposal level and does not include a locality-preserving fidelity term.

\subsection{Contribution}\label{sec:intro-contribution}

We introduce a kernel-aware local-refinement framework for ReCom-style
MCMC. The core object is a single Laplacian energy on the merged-pair
subgraph $H = G[V_i \cup V_j]$:
\begin{equation}\label{eq:core-energy}
  E(x) \;=\; \alpha\, x^\top L\, x \;+\; \beta\, \|x - x_0\|^2,
  \qquad
  (1 - \varepsilon)\bar p \;\le\; p^\top x \;\le\; (1+\varepsilon)\bar p,
\end{equation}
where $L$ is the kernel-weighted graph Laplacian of $H$, $x_0$ encodes
the current bipartition, and $\bar p$ is the target district
population. The Laplacian term $x^\top L x$ counts (or approximates)
cut edges between the two sides; the fidelity term $\|x - x_0\|^2$
penalises deviation from the current assignment, enforcing locality of
each move; and the population balance constraint guarantees
$\varepsilon$-balance. The framework has three primary contributions:

\begin{enumerate}
  \item \textbf{A kernel-aware local refinement} (\S\ref{sec:model}).
    The Laplacian $L$ accepts any nonnegative edge-weighting kernel.
    We instantiate the framework with three kernels --- uniform,
    perimeter-weighted, and density-weighted --- and demonstrate that
    each yields qualitatively different splits on Iowa
    (\S\ref{sec:exp-multikernel}). No prior MCMC redistricting sampler
    accepts arbitrary kernel weights at the proposal level.

  \item \textbf{A unified continuous/integer dual} (\S\ref{sec:model},
    \S\ref{sec:iqp}). The same energy admits two solvers:
    \begin{itemize}
      \item \textbf{Chain B (QP-diffusion)}: $x \in [0,1]^{|V_H|}$,
        convex QP, Cheeger-sweep rounding with orphan-flip repair.
      \item \textbf{Chain C (IQP local)}: $x \in \{0,1\}^{|V_H|}$,
        mixed-integer QP solved by Gurobi, orphan-flip repair on the
        integer optimum.
    \end{itemize}
    Both chains mix with vanilla or tempered ReCom for ergodicity.
    Setting $\alpha = 1, \beta = 0$, and the uniform kernel recovers
    the discrete cut-minimization integer program as a special case of
    Chain C; setting $\beta = 0$ and the uniform kernel recovers
    SpecReCom \cite{davies2025spectral} as a special case of Chain B.

\end{enumerate}

\noindent We also provide two supporting contributions that motivate
and validate the framework:

\begin{enumerate}\setcounter{enumi}{2}
  \item \textbf{A structural diagnostic of the partition fiber}
    (\S\ref{sec:fiber-corridors}). On three fully enumerated test beds
    (5$\times$5 at $\varepsilon = 0$, $\varepsilon = 0.20$; 4$\times$4
    at $\varepsilon = 0.25$), the partition fiber decomposes into
    \emph{islands} connected by \emph{narrow corridors} or
    \emph{single-edge bridges}. The narrow-throat count is
    approximately invariant ($\sim 12$--16) across grid size and
    balance tolerance. A chain step-trace diagnostic shows that
    ReCom jumps between islands at $\sim 13$--14\% per step while
    spectral local kernels stay within at $\sim 5$--7\% per step ---
    direct empirical confirmation that ReCom and the diffusion
    proposals address complementary regions of state space.

  \item \textbf{Empirical validation} (\S\ref{sec:experiments}). On
    Iowa precincts ($k = 4$, $\varepsilon = 0.05$), Chain C with
    $n_{\mathrm{iqp}} = 10$ delivers a 23.9\% reduction in mean cut
    edges versus Tempered ReCom ($p < 10^{-197}$), and Chain B
    delivers 3.2\% ($p = 3 \times 10^{-5}$). On the fully enumerated
    $5 \times 5$ ground truth, the total-variation distance to the
    uniform-on-fiber target distribution is bounded and characterized,
    quantifying the compactness-coverage trade-off.
\end{enumerate}

\paragraph{Positioning.} The framework relates to recent work as
follows. \emph{vs.\ SpecReCom}~\cite{davies2025spectral}: ours adds a
fidelity term that keeps moves local and unlocks kernel flexibility,
plus an integer dual. \emph{vs.\ RevReCom}~\cite{cannon2022spanning}:
ours mixes with reversible or tempered ReCom rather than replacing
the spanning-tree step, so reversibility against an explicit target
is preserved. \emph{vs.\ Cycle Walk, Marked Edge Walk,
BUD}~\cite{deford2025cycle, mcwhorter2025marked, akitaya2026bud}:
those chains operate at the spanning-forest level and target
mixing-time bounds; ours operates one level above on the merged-pair
subgraph, so the two regimes are complementary.

\paragraph{Roadmap.}\label{sec:intro-roadmap}
\S\ref{sec:related} reviews related work. \S\ref{sec:background} sets
up notation, recalls ReCom, and characterizes the partition fiber's
island-and-corridor macro-structure on three test beds.
\S\ref{sec:model} formulates the kernel-flexible QP-diffusion model in
both continuous (Chain B) and integer (Chain C) forms, with
Cheeger-sweep rounding, orphan-flip repair, hybrid composition with
ReCom, and a brief spectral interpretation.
\S\ref{sec:experiments} reports empirical results on Iowa, $5 \times 5$,
and \draft{NC precinct (pending Sprint 3)}. \S\ref{sec:discussion}
discusses trade-offs and limitations. Appendices contain algorithm
details, hyperparameter tables, full sweep tables, and a
reproducibility checklist.

%════════════════════════════════════════════════════════════════════════════════
\section{Related work}\label{sec:related}
%════════════════════════════════════════════════════════════════════════════════

\paragraph{MCMC for redistricting.}\label{sec:related-mcmc}

Early MCMC for redistricting framed the problem as a graph cut and used
single-precinct flip chains~\cite{fifield2020automated} or
Swendsen--Wang-style large-swap moves. Flip chains mix poorly and
produce highly irregular districts~\cite{deford2021recombination,
najt2024empirical}. ReCom~\cite{deford2021recombination} introduced
the spanning-tree merge-split move and is now standard. Extensions
include Multi-Scale Merge--Split~\cite{autry2020multiscale} preserving
nested communities (counties, precincts), Forest ReCom~\cite{autry2023metropolized},
Sequential Monte Carlo~\cite{mccartan2023sequential}, and reversible
recombination~\cite{cannon2022spanning} targeting the canonical
spanning-tree null distribution. Mixing-time bounds remain open for
all variants in the practical regime; Akitaya et
al.~\cite{akitaya2022reconfiguration} prove that even deciding
reachability between balanced partitions under recombination is
PSPACE-complete.

\paragraph{Spectral graph partitioning for redistricting.}\label{sec:related-spectral}

Spectral graph partitioning has a long history in image segmentation
and clustering~\cite{shi2000normalized}. In redistricting, the recent
SpecReCom of Davies et al.~\cite{davies2025spectral} replaces the
random spanning tree with a Fiedler-vector spectral cut, with
randomized edge weights to break determinism. SpecReCom empirically
yields more compact partitions than vanilla ReCom on grids and planar
maps. We discuss the relationship to our framework in
\S\ref{sec:intro-contribution} and \S\ref{sec:iqp-subsumption}.

\paragraph{Image segmentation and diffuse interface models.}\label{sec:related-image}

Diffuse interface models in Euclidean space are typically built around
the Ginzburg--Landau (GL) functional, which combines a smoothness term
$\int_\Omega |\nabla x|^2$ with a fidelity or potential term. The
Chan--Vese active contour model~\cite{chan2001active} and the
Merriman--Bence--Osher (MBO) threshold dynamics scheme implement this
energy iteratively, alternating smoothing with thresholding. On
graphs, the smoothness term becomes the Laplacian quadratic form
$\sum_{(i,j) \in E} w_{ij}(x_i - x_j)^2 = x^\top L x$, and graph
Tikhonov regularisation~\cite{pilavci2021graph,shuman2013emerging}
recovers the same heat-kernel-based smoothing. Bertozzi and
collaborators~\cite{bertozzi2012diffuse} apply the GL energy directly
to graph-based classification; the present work is the first to apply
this lineage to redistricting MCMC.

\paragraph{Cut minimization and discrete compactness.}\label{sec:related-cut}

The minimum graph cut~\cite{stoer1997simple,karger2000minimum} and its
balanced variants such as normalized cut~\cite{shi2000normalized} are
classical objectives. NP-hardness of balanced minimum
cut~\cite{wagner1993between} motivates continuous relaxations, of
which the Fiedler-vector cut and the Total Variation Isoperimetric
Profile~\cite{deford2019total} are two well-known examples. The cut
edge count $|\partial P|$ has been proposed as a discrete compactness
measure for redistricting~\cite{duchin2018discrete} and shown to
correlate well with classical geometric scores (Polsby--Popper, Reock)
while being more stable under unit refinement.

%════════════════════════════════════════════════════════════════════════════════
\section{Background: ReCom and the partition fiber}\label{sec:background}
%════════════════════════════════════════════════════════════════════════════════

\subsection{Dual graph and \texorpdfstring{$k$}{k}-district plans}\label{sec:background-graph}

Let $G = (V, E)$ be the dual graph of a state, with vertices
representing geographic units (precincts, census blocks, counties) and
edges encoding shared boundaries. Each vertex $v$ carries a population
$p_v \ge 0$ and may carry additional attributes (demographic vectors,
shapefile geometry). A \emph{$k$-district plan} is a partition
$P = \{V_1, \ldots, V_k\}$ of $V$ such that each induced subgraph
$G[V_i]$ is connected (\emph{contiguity}) and the populations satisfy
\begin{equation}\label{eq:pop-balance}
  (1 - \varepsilon)\bar p \;\le\; p(V_i) \;\le\; (1 + \varepsilon)\bar p,
  \qquad \bar p := \frac{p(V)}{k},
\end{equation}
for a tolerance $\varepsilon \in [0, 1)$. We write $\mathcal{P}_k(G; \varepsilon)$
for the set of valid plans. The \emph{cut edges} of a plan are
\[
  \partial P \;=\; \{(u,v) \in E : P(u) \ne P(v)\},
\]
and we call $|\partial P|$ the \emph{cut-edge count} of $P$. This
quantity is the discrete analogue of the perimeter of the district
boundaries and serves as our primary compactness
metric~\cite{duchin2018discrete}.

\begin{figure}[H]
  \centering
  \IfFileExists{fig/redistricting.png}{%
    \includegraphics[width=0.7\textwidth]{fig/redistricting.png}%
  }{%
    \fbox{\begin{minipage}{0.7\textwidth}\vspace{2.5cm}\centering
      \ttfamily\small\textcolor{orange}{[Missing figure:
      \texttt{fig/redistricting.png}]}\vspace{2.5cm}\end{minipage}}%
  }
  \caption{The dual graph of Iowa's 99 counties (left) and the current
    four congressional districts (right). Adapted
    from~\cite{deford2021recombination}.}
  \label{fig:iowa-dual}
\end{figure}

\subsection{Recombination (ReCom)}\label{sec:background-recom}

Recombination~\cite{deford2021recombination} is a Markov chain on
$\mathcal{P}_k(G; \varepsilon)$ whose single step is:

\begin{algorithm}[H]
\caption{Vanilla ReCom step (Chain V)}\label{alg:recom}
\KwIn{Current plan $P$, dual graph $G$, balance tolerance $\varepsilon$.}
\KwOut{Next plan $P'$.}
Select two adjacent districts $V_i, V_j$ uniformly at random from boundary pairs;\\
Form the induced subgraph $H = G[V_i \cup V_j]$;\\
\Repeat{a cut $e^*$ exists yielding two $\varepsilon$-balanced components}{
  Sample a uniform random spanning tree $T$ of $H$ via Wilson's algorithm~\cite{wilson1996generating};\\
  Search edges of $T$ for a cut $e^*$ whose removal produces two $\varepsilon$-balanced components.
}
Define $V_i', V_j'$ as the two components of $T \setminus \{e^*\}$;\\
$P' \gets P$ with $V_i, V_j$ replaced by $V_i', V_j'$;\\
\Return $P'$.
\end{algorithm}

\noindent The tree-based proposal flips many precincts simultaneously
between $V_i$ and $V_j$ while preserving contiguity automatically (any
edge removal from a tree produces two connected components). This is
the key advantage over single-precinct flip chains, which traverse the
state space too slowly to be practically useful.

\textbf{Variants used as baselines.} We compare against five chains:

\begin{table}[H]
\centering
\small
\begin{tabular}{lll}
\toprule
Chain & Method & Mechanism \\
\midrule
V & Vanilla ReCom~\cite{deford2021recombination} & spanning-tree split, always accept \\
A & Tempered ReCom & ReCom + Metropolis--Hastings on cuts \\
B & ReCom + QP-diffusion mix (this work) & continuous QP, Cheeger sweep, orphan-flip repair \\
C & Tempered ReCom + IQP local (this work) & integer QP, MIQP, orphan-flip repair \\
D & ReCom + SpecReCom mix~\cite{davies2025spectral} & Fiedler + Cheeger sweep, no $\beta$, no repair \\
\bottomrule
\end{tabular}
\caption{The five chains compared in this paper.}
\label{tab:five-chains}
\end{table}

Tempered ReCom uses a Metropolis--Hastings filter targeting the
distribution $\pi(P) \propto \exp(-\lambda |\partial P|)$ with the
symmetric-$q$ approximation; see Appendix~\ref{app:hyperparams} for
the locked-in $\lambda = 0.10$.

\paragraph{ReCom's two weaknesses.}\label{sec:background-weakness}
ReCom is empirically strong but has two known weaknesses that motivate
the present work.

\textbf{(W1) Non-compact splits.} Uniform spanning trees of the
merged subgraph favor topologically irregular partitions because the
number of irregular subtrees is combinatorially much larger than the
number of compact ones~\cite{deford2021recombination,deford2019total}.
The cut edge is chosen \emph{after} the tree is drawn, so the
proposal cannot anticipate compactness.

\textbf{(W2) Kernel-blindness.} The uniform spanning tree weights
every edge equally. Shared-boundary length, population density,
demographic similarity, and administrative boundary attributes are
ignored at the proposal level. They can only enter via post-hoc
filtering, which raises rejection rates and is hard to tune.

The diffusion-based proposal of \S\ref{sec:model} addresses both
weaknesses by replacing the spanning-tree-cut step with a single
convex optimization that minimizes a kernel-weighted Laplacian energy
with a locality-preserving fidelity term. The energy
$x^\top L\, x = \sum_{(u,v) \in E} w_{uv} (x_u - x_v)^2$ equals the
weighted cut count for binary $x$, so minimizing it directly minimizes
the discrete perimeter.

\subsection{Partition fiber structure: islands and corridors}\label{sec:fiber-corridors}

This subsection documents a structural diagnostic of the partition fiber
$\mathcal{P}_k(G; \varepsilon)$, viewed as a graph under the
single-flip adjacency relation. The diagnostic motivates the
local-refinement design of \S\ref{sec:model} and is, we believe, of
independent interest for analyzing any redistricting MCMC.

\paragraph{Three test beds.}\label{sec:fiber-testbeds}

We construct three fully enumerated partition fibers spanning grid
size, district count, and balance tolerance:

\begin{table}[H]
\centering
\small
\begin{tabular}{lrrrrr}
\toprule
Fiber & Nodes & Edges & $k$ & $\varepsilon$ & Allowed district sizes \\
\midrule
$5 \times 5$, $\varepsilon = 0$    & 4\,006   & 13\,416  & 5 & 0    & \{5\} \\
$5 \times 5$, $\varepsilon = 0.20$ & 193\,128 & 969\,760 & 5 & 0.20 & \{4, 5, 6\} \\
$4 \times 4$, $\varepsilon = 0.25$ & 1\,953   & 7\,024   & 4 & 0.25 & \{3, 4, 5\} \\
\bottomrule
\end{tabular}
\caption{The three test-bed fibers. Adjacency is single-flip: two
  partitions are connected by an edge iff one is obtained from the
  other by reassigning a single vertex.}\label{tab:fibers}
\end{table}

For each fiber, all valid partitions are enumerated and the single-flip
adjacency graph is built explicitly. \draft{Implementation:
\texttt{diffusion\_model/recom/fiber\_5x5/build\_fiber\_*.py};
runtimes on a 2024 MacBook range from $\sim 5$ s ($4 \times 4$,
$\varepsilon = 0.25$) to $\sim 50$ min ($5 \times 5$, $\varepsilon = 0.20$).}

\paragraph{Island structure: Louvain community decomposition.}\label{sec:fiber-islands}

Each fiber graph admits a clean community decomposition under the
Louvain modularity-maximization algorithm:

\begin{table}[H]
\centering
\small
\begin{tabular}{lrrr}
\toprule
Fiber & Communities & Modularity $Q$ & Community size range \\
\midrule
$5 \times 5$, $\varepsilon = 0$    & 10 & 0.65 & \needdata{compute} \\
$5 \times 5$, $\varepsilon = 0.20$ & 26 & 0.69 & $\sim 3\,000$--$13\,350$ \\
$4 \times 4$, $\varepsilon = 0.25$ & 16 & 0.61 & 59--170 \\
\bottomrule
\end{tabular}
\caption{Community decompositions. Modularity values $Q \ge 0.6$
  indicate strong community structure across all three fibers.}
\label{tab:communities}
\end{table}

\begin{figure}[H]
  \centering
  \begin{subfigure}[t]{0.48\textwidth}
    \centering
    \needdata{include fiber\_5x5/fiber\_5x5\_eps02\_corridors.png}\\
    \caption{$5 \times 5$ at $\varepsilon = 0.20$.}
    \label{fig:fiber-corridor-5x5}
  \end{subfigure}\hfill
  \begin{subfigure}[t]{0.48\textwidth}
    \centering
    \needdata{include fiber\_5x5/fiber\_4x4\_eps025\_corridors.png}\\
    \caption{$4 \times 4$ at $\varepsilon = 0.25$.}
    \label{fig:fiber-corridor-4x4}
  \end{subfigure}
  \caption{Two-level visualisation of the partition fiber:
    community centroids by spring layout on the meta-graph, nodes
    within each community by spring layout on the induced subgraph.
    Crimson edges indicate top-betweenness corridors (Brandes,
    $k = 100$ pivots).}
  \label{fig:fiber-corridors}
\end{figure}

\paragraph{Bridges, corridors, and the narrow-throat invariant.}\label{sec:fiber-corridors-invariant}

The salient feature of these decompositions is the connectivity pattern
between communities. At $\varepsilon = 0$ the fiber graph contains
\emph{single-edge bridges}; at $\varepsilon > 0$ these bridges thicken
into \emph{narrow corridors} of 2--4 edges, but their count is
preserved.

\begin{table}[H]
\centering
\small
\begin{tabular}{lrrrr}
\toprule
Fiber & Communities & Bridges (single-edge) & Narrow throats & Threshold \\
\midrule
$5 \times 5$, $\varepsilon = 0$    & 10 & \textbf{16} & ---          & ---         \\
$5 \times 5$, $\varepsilon = 0.20$ & 26 & 0           & \textbf{16}  & $<$50 edges \\
$4 \times 4$, $\varepsilon = 0.25$ & 16 & 0           & \textbf{13}  & $<$5 edges  \\
\bottomrule
\end{tabular}
\caption{The narrow-throat invariant. The number of low-traffic
  connection points between large communities lies in the order of
  a dozen across both grid sizes and balance tolerances.}
\label{tab:narrow-throats}
\end{table}

\paragraph{Provisional thesis.} Narrow-throat count is approximately a
topological invariant of the partition fiber's macro-structure across
$(\text{grid}, k, \varepsilon)$ variations. As $\varepsilon$ relaxes,
single-edge bridges thicken into few-edge corridors but the
cluster-connectivity skeleton retains its topology. \draft{A
theoretical explanation is open; we conjecture it relates to the
small number of pivot-precinct identities needed to transform one
balanced partition into another.}

\paragraph{Chain step-trace diagnostic.}\label{sec:fiber-trace}

The island/corridor decomposition supports a direct diagnostic for any
redistricting chain. We label each visited partition by its Louvain
community and record the \emph{jump rate} --- the fraction of steps
that move to a different community.

\begin{table}[H]
\centering
\small
\begin{tabular}{lrrr}
\toprule
Chain & Communities visited & Jump rate & Jumps in 1\,000 steps \\
\midrule
V: Vanilla ReCom       & 7/10 & \textbf{14.1\%} & 141 \\
A: Tempered ReCom      & 8/10 & \textbf{13.3\%} & 133 \\
B: QP-diffusion (50/50) & 6/10 & \textbf{5.2\%}  & 52  \\
C: ReCom + IQP local   & 5/10 & 11.5\%           & 115 \\
D: ReCom + SpecReCom   & 7/10 & \textbf{6.8\%}  & 68  \\
\bottomrule
\end{tabular}
\caption{Inter-community jump rates on the $5 \times 5$,
  $\varepsilon = 0$ fiber (1\,000 steps per chain). ReCom variants
  (V, A) jump between islands at $\sim 14\%$ per step; spectral local
  kernels (B, D) drop the jump rate to $\sim 5$--7\%.}\label{tab:island-trace}
\end{table}

\begin{figure}[H]
  \centering
  \needdata{include iowa\_hybrid\_test\_out/island\_trace\_5x5.png}\\
  \caption{Chain step-traces over $5 \times 5$ communities. Each panel
    shows the sequence of community labels visited by one chain over
    1\,000 steps; horizontal bands are communities, x-axis is time.
    Vanilla and Tempered ReCom oscillate rapidly between communities
    (high jump rate); QP-diffusion and SpecReCom dwell within
    communities for extended runs.}
  \label{fig:island-trace}
\end{figure}

\textbf{Direct empirical confirmation.} ReCom is for \emph{inter-island}
exploration; spectral local kernels are for \emph{intra-island}
exploration. The two regimes are complementary, motivating the hybrid
chain design of \S\ref{sec:model}.

\paragraph{Implication for chain design.}\label{sec:fiber-implication}

The diagnostic suggests an explicit principle:

\begin{quote}
\itshape A good redistricting sampler should combine global moves that
traverse narrow throats between islands with local moves that explore
each island's interior. ReCom does the former; spectral local kernels
do the latter.
\end{quote}

This is the design principle of Chains B and C: each step is a random
mixture of a ReCom (inter-island) move and a diffusion (intra-island)
move. The remainder of the paper develops and validates this design.

%════════════════════════════════════════════════════════════════════════════════
\section{The QP-diffusion model: continuous and integer formulations}\label{sec:model}
%════════════════════════════════════════════════════════════════════════════════

\subsection{Locality kernels}\label{sec:kernels}

Choosing the edge weights $w_{ij}$ is the most consequential modeling
decision. The kernel determines what \emph{smoothness} means and
therefore controls how district boundaries move under the relaxation.
We restrict all weights to geographic neighbors,
\begin{equation}\label{eq:kernel-local}
  w_{ij} = 0 \quad \text{if } (v_i, v_j) \notin E,
\end{equation}
ensuring that diffusion operates only through geographic adjacency and
preventing long-range interactions that could produce disconnected
districts.

\paragraph{Uniform kernel ($k_{\mathrm{unif}}$).} $w_{ij} = 1$ for
every edge. The Laplacian $L$ is the combinatorial graph Laplacian;
$x^\top L x$ is exactly the cut edge count. Recovers SpecReCom
under $(\beta = 0, k_{\mathrm{unif}})$.

\paragraph{Perimeter kernel ($k_{\mathrm{perim}}$).} $w_{ij} = |\partial v_i \cap \partial v_j|$
(shared-boundary length). The Laplacian quadratic form $x^\top L x$
is then the \emph{geographic perimeter} of the partition, weighted by
boundary segment length. This kernel discourages fragmented districts
and favors boundaries that cut along short segments.

\paragraph{Density kernel ($k_{\mathrm{dens}}$).} Let $\rho_v$ denote
the population density at vertex $v$ (population per unit area). A
Gaussian density kernel
\[
  w_{ij} \;=\; |\partial v_i \cap \partial v_j| \cdot
  \exp\!\Bigl( -\frac{(\rho_i - \rho_j)^2}{\sigma^2} \Bigr)
\]
encourages cuts along sharp density gradients (urban--rural
boundaries) while preserving the topological role of shared boundary
length.

\paragraph{Demographic kernel ($k_{\mathrm{demog}}$).} Let
$d_v \in \R^m$ be a vector of demographic attributes. A Gaussian
similarity kernel $w_{ij} = \exp(-\|d_i - d_j\|^2 / \sigma^2)$
encourages communities of interest. Should be used as a soft
regulariser, not a hard constraint, to avoid demographic engineering.

\paragraph{Composite kernels.} Linear combinations $w_{ij} = \sum_k a_k w_{ij}^{(k)}$
with $a_k \ge 0$ remain valid (preserve nonnegativity and locality).

\subsection{The convex QP}\label{sec:qp-model}

Let $H = G[V_i \cup V_j]$ be the merged subgraph of two adjacent
districts. Let $L_H$ be the kernel-weighted Laplacian of $H$, let
$x_0 \in \{0,1\}^{|V_H|}$ encode the current bipartition (1 for
vertices in $V_i$, 0 for vertices in $V_j$), let $p$ be the population
vector restricted to $H$, and let $\bar p = p(V_H)/2$ be the target
population per side. Fix weights $\alpha, \beta > 0$ and balance
tolerance $\varepsilon \in (0, 1)$.

\begin{definition}[Continuous QP-diffusion]\label{def:qp}
The continuous QP-diffusion bipartition of $H$ is
\begin{align}
  x^\star \;=\; \arg\min_{x \in [0,1]^{|V_H|}} \;\;
    & \alpha\, x^\top L_H\, x \;+\; \beta\, \|x - x_0\|_2^2 \label{eq:qp-obj} \\
  \text{subject to} \;\;
    & (1 - \varepsilon) \bar p \;\le\; p^\top x \;\le\; (1 + \varepsilon) \bar p.
    \label{eq:qp-balance}
\end{align}
\end{definition}

The first term in~\eqref{eq:qp-obj} is the kernel-weighted Laplacian
smoothness energy: for binary $x$ it equals the weighted cut count
$\sum_{(u,v) \in E(H)} w_{uv}\,\mathbf{1}[x_u \ne x_v]$, and the
continuous relaxation $x \in [0,1]^{|V_H|}$ extends this to a smooth
objective on the unit cube. The second term is the \emph{fidelity term},
which penalises deviation from the current bipartition and ensures
that each proposed move is local. The balance constraint enforces
$\varepsilon$-balance on the proposed split.

\begin{proposition}[Convexity and well-posedness]\label{prop:convex}
The objective in~\eqref{eq:qp-obj} is strictly convex on
$[0,1]^{|V_H|}$ (since $\beta > 0$ and $L_H$ is positive semidefinite).
The feasible region defined by~\eqref{eq:qp-balance} is a non-empty
polytope (it contains $x_0$). Hence the QP admits a unique minimizer
$x^\star$, computable in polynomial time by any convex QP solver.
\end{proposition}

\paragraph{Spectral interpretation.} Projecting the interior optimality
condition $(\alpha L_H + \beta I) x^\star = \beta x_0$ onto the
eigenbasis of $L_H$ with eigenvalues
$0 = \lambda_1 \le \lambda_2 \le \cdots \le \lambda_n$ yields the
closed-form frequency response
\begin{equation}\label{eq:freq-response}
  \hat x^\star_k \;=\; \frac{\beta}{\alpha \lambda_k + \beta}\,
  \hat x_{0,k},
  \qquad \hat y_k := \phi_k^\top y.
\end{equation}
The QP therefore acts as a low-pass spectral filter on the prior $x_0$,
with cutoff controlled by the ratio $\alpha / \beta$. Equivalently, with
$\tau := \alpha / \beta$, the optimum satisfies
$(x^\star - x_0)/\tau = -L_H x^\star$ --- the implicit Euler
discretisation of the graph heat equation $\dot x = -L_H x$. Inspecting
the leading eigenvectors $\phi_2, \phi_3, \ldots$ provides a sanity
check on kernel choice: if they look like noise rather than coherent
geographic regions, the chosen kernel is probably not informative.

\subsection{Cheeger-sweep rounding and orphan-flip repair}\label{sec:rounding}

The continuous optimum $x^\star \in [0,1]^{|V_H|}$ must be rounded to a
binary bipartition. We use a \emph{Cheeger sweep}~\cite{kelner2009spectral}:
sort the entries of $x^\star$ in descending order, and for each
threshold $k = 1, \ldots, |V_H| - 1$ form the candidate set $S_k$ of the
top-$k$ vertices. Among candidates that satisfy $\varepsilon$-balance,
return the one minimizing the cut count
$|\partial(S_k, V_H \setminus S_k)|$. Candidates with disconnected sides
are passed through the orphan-flip repair below before evaluation.

\begin{algorithm}[H]
\caption{Cheeger-sweep rounding with orphan-flip repair}\label{alg:cheeger-sweep}
\KwIn{Continuous optimum $x^\star \in [0,1]^{|V_H|}$, population
  vector $p$, target $\bar p$, tolerance $\varepsilon$, merged subgraph $H$.}
\KwOut{Binary bipartition $S \subset V_H$ or \texttt{FAIL}.}
Sort $V_H$ by $x^\star$ descending; call the order $v_{(1)}, \ldots, v_{(|V_H|)}$;\\
$\mathcal{C} \gets \emptyset$ \tcp*{candidates}
\For{$k = 1$ \KwTo $|V_H| - 1$}{
  $S_k \gets \{v_{(1)}, \ldots, v_{(k)}\}$;\\
  \uIf{$p(S_k) \in [(1-\varepsilon)\bar p, (1+\varepsilon)\bar p]$ and both sides connected}{
    Add $S_k$ to $\mathcal{C}$ with cut $|\partial(S_k, V_H \setminus S_k)|$;
  }\ElseIf{$p(S_k)$ in balance and sides have $> 1$ components}{
    $S_k' \gets$ orphan-flip repair of $S_k$ (Algorithm~\ref{alg:orphan-flip});\\
    \If{$S_k'$ satisfies balance and connectivity}{
      Add $S_k'$ to $\mathcal{C}$;
    }
  }
}
\lIf{$\mathcal{C} = \emptyset$}{\Return \texttt{FAIL}}
\Return $\arg\min_{S \in \mathcal{C}} |\partial(S, V_H \setminus S)|$.
\end{algorithm}

The sweep can return a candidate $S_k$ whose induced subgraph $G[S_k]$
or $G[V_H \setminus S_k]$ is disconnected --- a single boundary vertex
on the wrong side of the threshold is enough. Discarding such
candidates outright would inflate the rejection rate; instead we repair
them. The \emph{orphan-flip} routine identifies the disconnected
components of each side, keeps the largest component, and flips every
smaller (\emph{orphan}) component to the opposite side. Population
balance is then re-verified; if the repair preserves both balance and
connectivity, the repaired candidate enters the Cheeger sweep's
candidate pool. Algorithm~\ref{alg:orphan-flip} formalizes this:

\begin{algorithm}[H]
\caption{Orphan-flip repair}\label{alg:orphan-flip}
\KwIn{Candidate bipartition $(S, V_H \setminus S)$.}
\KwOut{Repaired bipartition or \texttt{FAIL}.}
Identify connected components of $G[S]$ and $G[V_H \setminus S]$;\\
\For{each side}{
  Keep the largest connected component;\\
  Flip every smaller component to the opposite side;
}
Verify both final sides are connected and $\varepsilon$-balanced;\\
\lIf{verification fails}{\Return \texttt{FAIL}}
\Return repaired $(S', V_H \setminus S')$.
\end{algorithm}

\textbf{Why repair matters.} On heterogeneous-population graphs such
as Iowa, the unrepaired Cheeger sweep produces near-feasible
candidates whose few orphan vertices --- often singletons at the
boundary --- can be reassigned without affecting the cut count.
Without repair, post-rounding rejection rates on Iowa run at 5--15\%;
with repair, $<3\%$. SpecReCom~\cite{davies2025spectral} does not use
repair (and we keep it out of our faithful Chain D
implementation); without the $\beta$ fidelity term, repair tends to
accept lower-quality cuts.

\paragraph{Cheeger bound on the rounded cut.} By the Cheeger
inequality~\cite{kelner2009spectral}, the Cheeger constant of $H$ is
bounded by $\Phi(H) \le \sqrt{2 \lambda_2(L_H^{\mathrm{sym}})}$, where
$\lambda_2$ is the second eigenvalue of the normalized Laplacian. The
sweep applied to $x^\star$ inherits this bound on the
isoperimetric ratio of $S^\star$, since the leading-frequency content
of $x^\star$ in~\eqref{eq:freq-response} dominates by the
low-pass-filter property.

\subsection{Hybrid composition with ReCom}\label{sec:hybrid-mh}

Chain B's single step is a random mixture: with probability
$p_R = 0.5$, take a Vanilla ReCom step (Algorithm~\ref{alg:recom});
otherwise take a QP-diffusion step (Definition~\ref{def:qp} +
Algorithm~\ref{alg:cheeger-sweep}).

\begin{proposition}[Ergodicity of the hybrid]\label{prop:ergodic}
If ReCom is ergodic on $\mathcal{P}_k(G; \varepsilon)$ (which holds
under standard assumptions; see~\cite{deford2021recombination,
cannon2024irreducibility}), then the convex combination of ReCom and
the QP-diffusion proposal with any $p_R \in (0, 1)$ is also ergodic.
\end{proposition}

\begin{proof}[Proof sketch]
Irreducibility of the mixture follows from irreducibility of ReCom
alone: any plan reachable by ReCom in $T$ steps remains reachable in
expected $T / p_R$ mixture steps. Aperiodicity follows because the
QP-diffusion step has positive probability of returning the current
plan (when the convex optimum is already binary), giving the mixture
chain a self-loop probability.
\end{proof}

The Metropolis--Hastings filter for the diffusion step targets the
tempered distribution $\pi(P) \propto \exp(-\lambda |\partial P|)$ with
$\lambda = 0.10$, using the symmetric-$q$ approximation:
\[
  \alpha_{\mathrm{MH}}(P \to P') \;=\; \min\!\bigl(1,\,
    \exp(-\lambda(|\partial P'| - |\partial P|))\bigr).
\]
The symmetric-$q$ approximation introduces a small bias on small
graphs (see \S\ref{sec:exp-5x5}); on Iowa the bias is negligible
relative to the compactness gain.

\subsection{The integer QP-diffusion (Chain C)}\label{sec:iqp}

The integer counterpart of Definition~\ref{def:qp} replaces
$x \in [0,1]^{|V_H|}$ with $x \in \{0, 1\}^{|V_H|}$ and solves the
resulting mixed-integer convex quadratic program (MIQP) directly.

\paragraph{Definition and solver.}\label{sec:iqp-def}

\begin{definition}[Integer QP-diffusion]\label{def:iqp}
The IQP local refinement of $H$ is
\begin{align}
  x^\star \;=\; \arg\min_{x \in \{0,1\}^{|V_H|}} \;\;
    & \alpha\, x^\top L_H\, x \;+\; \beta\, \|x - x_0\|_2^2 \label{eq:iqp-obj} \\
  \text{subject to} \;\;
    & (1 - \varepsilon) \bar p \;\le\; p^\top x \;\le\; (1 + \varepsilon) \bar p.
\end{align}
\end{definition}

The objective is convex quadratic and the variable is binary.
Gurobi's MIQP solver handles this directly via branch-and-bound; in
practice, a 5-second time limit per solve is rarely hit on Iowa-sized
merged pairs ($\sim 30$--$60$ precincts).

\paragraph{Subsumption of cut-min and SpecReCom.}\label{sec:iqp-subsumption}

\begin{proposition}[Special cases]\label{prop:iqp-special-cases}
\begin{enumerate}
  \item Setting $\alpha = 1$, $\beta = 0$, $w_{ij} \equiv 1$ in
    Definition~\ref{def:iqp} gives the classical minimum-cut integer
    program subject to $\varepsilon$-balance. Hence Chain C subsumes
    the discrete cut-minimization MIP as a special case.
  \item Setting $\beta = 0$, $w_{ij} \equiv 1$ in
    Definition~\ref{def:qp} (continuous case) and rounding via the
    Cheeger sweep gives (up to weight randomization) the SpecReCom
    proposal of~\cite{davies2025spectral}. Hence Chain B subsumes
    SpecReCom as a special case.
\end{enumerate}
\end{proposition}

\paragraph{Locality via the fidelity term.}\label{sec:iqp-locality}

For $\beta = 0$, the IQP optimum is a minimum cut of $H$ subject only
to population balance and may differ arbitrarily from $x_0$. For
$\beta > 0$, the optimum is biased toward $x_0$: any flip from $x_0$
incurs a $\beta$-penalty per bit changed. This has three coupled
benefits over $\beta = 0$:

\begin{enumerate}
  \item \textbf{Locality.} The proposed bipartition stays geographically
    close to the current one, so the chain explores neighboring
    plans rather than randomly resetting on each merge.
  \item \textbf{Contiguity-friendliness.} Because $x_0$ encodes a
    contiguous bipartition and the IQP optimum is biased toward $x_0$,
    the optimum is mostly contiguous already; orphan-flip repair
    needs only to fix a small number of boundary-flipping vertices.
  \item \textbf{Lower rejection.} Post-rounding rejection on Iowa
    drops from 5--15\% ($\beta = 0$) to $<3\%$ ($\beta = 1$) with
    repair, while cut quality is preserved.
\end{enumerate}

\paragraph{One step of Chain C.}\label{sec:chain-c-step}
A single step of Chain C wraps an outer Tempered ReCom move
(Algorithm~\ref{alg:recom} plus the Metropolis--Hastings filter of
\S\ref{sec:hybrid-mh}) around $n_{\mathrm{iqp}}$ deterministic
IQP local-refinement passes. Each inner pass picks an adjacent district
pair, solves the IQP on the merged subgraph, and applies orphan-flip
repair; if the repair fails the inner pass leaves the plan unchanged.
The full step is given as Algorithm~\ref{alg:chain-c}.

\begin{algorithm}[H]
\caption{One step of Chain C}\label{alg:chain-c}
\KwIn{Current plan $P$, parameters $\alpha, \beta, n_{\mathrm{iqp}}, \lambda, \varepsilon$.}
\KwOut{Next plan $P'$.}
$P_0 \gets $ Tempered ReCom step (Chain A) applied to $P$\;
$P_{\mathrm{curr}} \gets P_0$\;
\For{$t = 1$ \KwTo $n_{\mathrm{iqp}}$}{
  Select adjacent district pair $(V_i, V_j)$ uniformly from boundary pairs of $P_{\mathrm{curr}}$\;
  Form $H = G[V_i \cup V_j]$, build kernel-weighted $L_H$\;
  Solve IQP (Definition~\ref{def:iqp}) for $x^\star$\;
  Apply orphan-flip repair (Algorithm~\ref{alg:orphan-flip})\;
  \uIf{repaired bipartition feasible and connected}{
    Update $V_i, V_j$ in $P_{\mathrm{curr}}$ to the new bipartition\;
  }\Else{
    $P_{\mathrm{curr}}$ unchanged\;
  }
}
\Return $P' \gets P_{\mathrm{curr}}$.
\end{algorithm}

The IQP loop is \emph{deterministic} given the random pair choice ---
no Metropolis acceptance step inside the loop. The MH acceptance is
inherited from the outer Tempered ReCom step; the IQP refinement is
treated as a within-step deterministic improvement against the same
target distribution. \todo{Verify that this construction preserves
reversibility against the tempered target. The standard argument is:
each IQP step is deterministic and so does not perturb the chain's
stationary distribution, provided the IQP map is detailed-balanced
with respect to $\pi$. This needs an explicit calculation; see
\S\ref{sec:discussion}.}

\paragraph{Pareto trade-off: continuous vs.\ integer.}\label{sec:iqp-pareto}

\begin{table}[H]
\centering
\small
\begin{tabular}{lll}
\toprule
 & Chain B (continuous QP) & Chain C (integer IQP) \\
\midrule
Variable & $x \in [0,1]^n$ & $x \in \{0,1\}^n$ \\
Solver & Convex QP (fast) & Convex MIQP (branch-and-bound) \\
Per-step cost & $O(\text{nnz}(L_H))$ & $O$(MIQP, $\le 5$s limit) \\
Cut quality (Iowa) & $-3.2\%$ vs A & $-23.9\%$ vs A \\
Contiguity & after rounding + repair & after repair (mostly already integer) \\
\bottomrule
\end{tabular}
\caption{Chain B and C trade-off. B is faster; C is more compact.
  Both kernel-flexible.}\label{tab:b-vs-c}
\end{table}

%════════════════════════════════════════════════════════════════════════════════
\section{Experiments}\label{sec:experiments}
%════════════════════════════════════════════════════════════════════════════════

We evaluate the five chains of Table~\ref{tab:five-chains} on three
datasets: the $5 \times 5$ grid (full enumeration, ground-truth
target), Iowa precincts (real geography, $k = 4$, $\varepsilon = 0.05$),
and \draft{NC precincts (k=4 sub-state; pending Sprint 3
completion)}. All hyperparameters are locked in
Appendix~\ref{app:hyperparams}; raw cut traces and chain outputs are
released in the project repository.

\subsection{Multi-kernel showcase on Iowa}\label{sec:exp-multikernel}

\needdata{This is the headline experiment per PAPER\_PLAN \S7.1. Run
Chains B and C with $k_{\mathrm{unif}}$, $k_{\mathrm{perim}}$,
$k_{\mathrm{dens}}$. Produce side-by-side Iowa partition maps.
Pending Sprint 3.}

We run Chain B and Chain C on Iowa with three kernels and visualise
the resulting partitions. The three kernels are uniform
($k_{\mathrm{unif}}$), perimeter ($k_{\mathrm{perim}}$,
$w_{ij}$ = shared-boundary length from \texttt{shared\_perim}), and
density ($k_{\mathrm{dens}}$, $w_{ij} = |\partial v_i \cap \partial v_j| \cdot \exp(-(\rho_i - \rho_j)^2 / \sigma^2)$
with $\rho_v = \mathrm{TOTPOP}/\mathrm{area}$).

\begin{figure}[H]
  \centering
  \needdata{three Iowa partition maps under $k_{\mathrm{unif}}$,
    $k_{\mathrm{perim}}$, $k_{\mathrm{dens}}$.}\\
  \caption{Iowa partitions under three kernels. Same merged-pair
    selections, same balance constraints; the kernel re-weights the
    Laplacian and shifts the chosen cut. The three plans differ
    visibly in which counties they group together, demonstrating that
    the kernel directly encodes geographic signal at the proposal
    level.}
  \label{fig:multikernel-iowa}
\end{figure}

\begin{table}[H]
\centering
\small
\begin{tabular}{lrrr}
\toprule
Kernel & Chain B mean cuts & Chain C mean cuts & Cuts vs uniform \\
\midrule
$k_{\mathrm{unif}}$  & 37.07 & 30.66 & ---           \\
$k_{\mathrm{perim}}$ & \needdata{} & \needdata{} & \needdata{} \\
$k_{\mathrm{dens}}$  & \needdata{} & \needdata{} & \needdata{} \\
\bottomrule
\end{tabular}
\caption{Cut counts under three kernels on Iowa ($k = 4$,
  $\varepsilon = 0.05$, 500 steps, burn-in 100).}\label{tab:multikernel-iowa}
\end{table}

\textbf{Novelty vs.\ SpecReCom.} SpecReCom uses unweighted spectral
splits with randomized edge weights; ReCom uses uniform spanning trees.
Neither encodes geographic signal at the proposal level. Our framework
is the first redistricting MCMC that accepts arbitrary kernel weights
at the proposal level, with both continuous and integer solvers.

\note{After the Sprint 3 multi-kernel runs, write a 2--3 paragraph
narrative interpreting what each kernel \emph{does to the Iowa
partition}, not just its cut count:
(1) does $k_{\mathrm{perim}}$ reduce the \emph{geographic} perimeter
(sum of shared boundary lengths in km) compared to $k_{\mathrm{unif}}$?
(2) does $k_{\mathrm{dens}}$ separate urban from rural counties at the
boundaries (compare to county-level density map)?
(3) which county pairs flip between districts under each kernel?
The headline of the paper is the partition maps and their
interpretation; the cut-count table is supporting data.}

\subsection{Compactness comparison across datasets}\label{sec:exp-compactness}

\paragraph{Iowa (5 chains).}\label{sec:exp-iowa}

\begin{table}[H]
\centering
\small
\begin{tabular}{lrrrr}
\toprule
Chain & Mean cuts & Std & $\Delta$ vs A & Wall-clock \\
\midrule
V: Vanilla ReCom         & 41.81 & 3.85 & ---                               & 0.2\,s \\
A: Tempered ReCom        & 38.29 & 3.02 & ---                               & 0.2\,s \\
B: QP-diffusion mix      & 37.07 & 4.14 & $-3.2\%$ ($p = 3 \times 10^{-5}$) & 1.1\,s \\
\textbf{C: IQP local ($n_{\mathrm{iqp}} = 10$)}
                         & \textbf{30.66} & \textbf{1.91}
                         & \textbf{$-23.9\%$ ($p < 10^{-197}$)}              & 9.9\,s \\
D: SpecReCom mix         & 40.79 & 4.47 & $+6.5\%$ (NS)                     & 0.6\,s \\
\bottomrule
\end{tabular}
\caption{Iowa cut counts (k = 4, $\varepsilon = 0.05$, 500 steps,
  burn-in 100). Chain C with $n_{\mathrm{iqp}} = 10$ delivers a 23.9\%
  reduction in mean cuts versus Tempered ReCom, with the tightest
  distribution (std 1.91). Chain B is an intermediate Pareto point.
  Chain D performs no better than Vanilla ReCom on Iowa --- the
  $\beta$ fidelity term is doing real work.}\label{tab:iowa-headline}
\end{table}

\begin{figure}[H]
  \centering
  \needdata{Pareto frontier plot: cuts vs wall-clock for V, A, B, C, D
    on Iowa. (Easy from existing npz outputs.)}\\
  \caption{Pareto frontier on Iowa: cut compactness against
    wall-clock time. Chain C dominates on cuts; Chain B is the
    intermediate point favored by practitioners with budget
    constraints.}
  \label{fig:pareto-iowa}
\end{figure}

\begin{figure}[H]
  \centering
  \needdata{include iowa\_hybrid\_test\_out/target\_matched\_iowa.png}\\
  \caption{Cut-edge trace and post-burn-in histogram for the five
    chains on Iowa. The diffusion-enhanced chains (B, C) reach
    low-cut configurations in fewer steps and concentrate around them
    tightly.}
  \label{fig:iowa-trace}
\end{figure}

\note{Add a narrative paragraph here interpreting
Table~\ref{tab:iowa-headline}:
(1) why does Chain D (SpecReCom) underperform Chain V here when both
strip out the $\beta$ term? Hypothesis: the $\beta$ fidelity term is
doing most of the work, not the spectral cut itself.
(2) Why does Chain C achieve a much tighter std (1.91) than B (4.14)?
(3) How does the wall-clock factor of $\sim 10\times$ between B and C
inform practitioner choice?
This is the empirical headline; the table cannot do the interpretation
on its own.}

\paragraph{5$\times$5 ground-truth validation and bias bound.}\label{sec:exp-5x5}

The $5 \times 5$ partition fiber ($k = 5$, $\varepsilon = 0$, 4\,006
partitions; \S\ref{sec:fiber-corridors}) admits full enumeration,
giving an exact uniform-on-fiber target distribution. We compare
empirical visit frequencies to this exact target.

\begin{table}[H]
\centering
\small
\begin{tabular}{lrrrr}
\toprule
Chain & Mean cuts & Std & Unique / 4006 & TV(emp, uniform) \\
\midrule
V: Vanilla ReCom     & 16.60 & 0.91 & 280 / 4006 (7.0\%) & 0.93 \\
A: Tempered ReCom    & 16.70 & 0.95 & 315 / 4006 (7.9\%) & 0.92 \\
B: QP-diffusion mix  & 16.67 & 0.95 & 164 / 4006 (4.1\%) & 0.96 \\
\textbf{C: IQP local}        & \textbf{16.17} & \textbf{0.47}
                              & 179 / 4006 (4.5\%) & 0.96 \\
D: SpecReCom mix     & 16.58 & 0.84 & 199 / 4006 (5.0\%) & 0.95 \\
\bottomrule
\end{tabular}
\caption{$5 \times 5$ ground-truth comparison ($k = 5$,
  $\varepsilon = 0$, 1\,000 steps, burn-in 100). Compactness gains
  trade against representativeness: V and A cover 7--8\% of the
  fiber; B and C cover 4--5\%. The trade-off is bounded and
  characterized.}\label{tab:5x5-headline}
\end{table}

\textbf{Honest cost of compactness.} The compactness-improving chains
(B, C) sacrifice approximately 3 percentage points of fiber coverage
versus V/A. The TV-to-uniform distance increases from $\sim 0.92$ to
$\sim 0.96$. This is the price of moving the empirical distribution
toward more compact partitions; the chain remains correctly
implemented against its target distribution (verified by the
diagnostic).

\note{Add 2--3 sentences interpreting why the trade-off is so much
\emph{smaller} on the 5$\times$5 grid than on Iowa
(3.2\% vs 23.9\% reduction): uniform vertex weights leave less
room for compactness improvement, whereas Iowa's heterogeneous
counties admit much larger gains. Argue that the 5$\times$5 result
is a \emph{calibration} of the framework, not a demonstration of its
power.}

\subsection{Ablations and sensitivity}\label{sec:exp-ablations}

\paragraph{Diversity controls: Jitter and Tabu.}\label{sec:exp-diversity}

\begin{table}[H]
\centering
\small
\begin{tabular}{lrrrr}
\toprule
Configuration & Mean cuts & Std & Coverage & Wall \\
\midrule
Baseline ($n_{\mathrm{iqp}} = 5$)               & 31.88 & 2.28 & 4.5\% & 39.7\,s \\
+ Jitter ($\mathcal{U}[1,2]$ on edge weights)   & 31.76 & 1.77 & 5.7\% & 45.2\,s \\
+ Tabu (size 3)                                  & \textbf{31.29} & \textbf{1.26} & 4.6\% & \textbf{37.6\,s} \\
\bottomrule
\end{tabular}
\caption{Diversity controls on Iowa ($n_{\mathrm{iqp}} = 5$). Tabu
  delivers the best combination of cut quality, distribution
  tightness, and runtime.}\label{tab:diversity-iowa}
\end{table}

Tabu, in particular, tightens the cut-count distribution
(std $2.28 \to 1.26$, a 44\% reduction) without sacrificing coverage
and at a modest runtime saving. Both controls are configurable knobs
for the practitioner trading compactness against representativeness.

\note{Add a paragraph explaining \emph{why} Jitter and Tabu help:
Jitter breaks ties between equally-good cuts, encouraging the chain
to explore equivalent configurations rather than locking onto one
deterministic optimum (which would happen with $\beta > 0$ and a tight
MIQP); Tabu forbids the same merged-pair from being picked
consecutively, forcing the chain to vary which boundary it refines.
Frame them as orthogonal: Jitter affects \emph{which cut} on a given
pair, Tabu affects \emph{which pair} the chain picks next.}

\paragraph{Parameter sensitivity: $n_{\mathrm{iqp}}$, $\alpha$, $\beta$.}\label{sec:exp-sensitivity}

\begin{table}[H]
\centering
\small
\begin{tabular}{lrrr}
\toprule
Configuration & Mean cuts & $\Delta$ vs A & Wall-clock \\
\midrule
$n = 3, \alpha = 10, \beta = 1$  & 32.37 & $-19.6\%$ & 24\,s \\
$n = 5, \alpha = 10, \beta = 1$  & 31.88 & $-20.8\%$ & 40\,s \\
\textbf{$n = 10, \alpha = 10, \beta = 1$} & \textbf{30.66} & \textbf{$-23.9\%$} & 103\,s \\
$n = 3, \alpha = 50, \beta = 1$  & 32.38 & $-19.6\%$ & 32\,s \\
$n = 3, \alpha = 10, \beta = 0$  & 32.31 & $-19.8\%$ & 31\,s \\
$n = 5, \alpha = 10, \beta = 0$  & 31.45 & $-21.9\%$ & 46\,s \\
\bottomrule
\end{tabular}
\caption{Chain C parameter sensitivity on Iowa. $n_{\mathrm{iqp}}$ is
  the dominant lever; $\alpha$ saturates; $\beta$ has minor effect
  on cut quality but governs the locality structure exploited by
  orphan-flip repair.}\label{tab:sensitivity}
\end{table}

\note{Add 1--2 paragraphs interpreting Table~\ref{tab:sensitivity}:
(1) Why does $n_{\mathrm{iqp}}$ scale roughly linearly with runtime
but \emph{sub-linearly} with cut reduction (i.e., diminishing returns)?
(2) Why does $\alpha$ saturate at 10? At what $\alpha$ does
oversmoothing begin?
(3) Why is $\beta = 0$ \emph{nearly as good} on cuts but with no
contiguity guarantee --- giving up locality buys almost nothing on
this metric but costs the rejection-rate advantage.}

\subsection{Pending experiments}\label{sec:exp-pending}

\paragraph{North Carolina precinct validation.}\label{sec:exp-nc}
\needdata{Pending Sprint 3 (PAPER\_PLAN \S7.5). NC precinct (2\,692
nodes, $k = 4$). Expect Chain C per-step cost $\sim 5$--30\,s; total
$\sim 10$--30 min for 100 steps. Cross-state replication of the
Iowa compactness gain.}

\paragraph{Downstream partisan-skew analysis.}\label{sec:exp-skew}
\needdata{Pending Sprint 4 (PAPER\_PLAN \S7.6). 5\,000 plans per chain
on Iowa. Compute efficiency gap, mean-median, packing. KS test for
distributional differences. Headline question: "does the compactness
improvement change the partisan answer?"}

%════════════════════════════════════════════════════════════════════════════════
\section{Discussion}\label{sec:discussion}
%════════════════════════════════════════════════════════════════════════════════

\note{This Discussion is currently too short (one paragraph plus bullet
lists). Expand to $\sim$2 pages of real synthesis:
(a) \emph{Why does the framework work?} What property of the
QP-diffusion explains the 23.9\% gain on Iowa?
(b) \emph{When should a practitioner choose B vs C?} Write explicit
guidance: B for $> 10^4$ steps on large states (wall-clock); C for
high-quality ensembles on $\le 10^3$ steps.
(c) \emph{Why does Chain D fail on Iowa?} Use this to argue that the
$\beta$ fidelity term is the load-bearing innovation, not the
spectral cut.
(d) \emph{What does the narrow-throat invariant mean?} A
two-paragraph reflection on whether this is a topological feature of
the redistricting problem itself or an artifact of the chosen
adjacency relation.
Once expanded, consider whether the three subsections below should
stay as subsections or be merged into 2--3 thematic paragraphs ---
short subsections in a short Discussion look like an unfinished
draft.}

\subsection{Compactness--coverage trade-off}\label{sec:disc-tradeoff}

The empirical results document a bounded and characterizable trade-off.
On the fully enumerated $5 \times 5$ ground truth
(Table~\ref{tab:5x5-headline}), Chain C achieves a 3.2\% reduction in
mean cuts compared to Tempered ReCom, at a cost of $\sim 3$ percentage
points in fiber coverage and a $0.04$ increase in TV distance to
uniform. The compactness gain on Iowa
(Table~\ref{tab:iowa-headline}) is much larger (23.9\%) because Iowa
admits real compactness improvements that the $5 \times 5$ uniform-population
grid does not. The Jitter and Tabu controls
(Table~\ref{tab:diversity-iowa}) give the practitioner explicit
configuration knobs for navigating this trade-off.

\subsection{Limitations}\label{sec:disc-limits}

\textbf{(L1) Continuous QP is loose on heterogeneous-population graphs.}
The continuous QP-diffusion (Chain B) gives only a 3.2\% cut reduction
on Iowa. The integer version (Chain C) recovers a 23.9\% reduction at
a $\sim 10\times$ compute cost. The continuous relaxation appears to
suffer on graphs with highly variable vertex populations; the integer
version is closer to optimal but slower.

\textbf{(L2) No theoretical mixing time bound.} Like every existing
redistricting chain, we lack a polynomial mixing-time guarantee.
We rely on empirical diagnostics: Gelman--Rubin $\hat R$, effective
sample size, and the island-trace diagnostic of
\S\ref{sec:fiber-trace}.

\textbf{(L3) Symmetric-$q$ approximation in Tempered ReCom.} The
Tempered ReCom Metropolis filter uses the symmetric-$q$ approximation,
which is biased on small graphs ($5 \times 5$: empirical mean cuts
16.70 vs.\ target 18.56). The bias is negligible on Iowa and other
large graphs.

\textbf{(L4) Continuous QP reversibility.} Chain B's continuous
QP-diffusion may not be inherently reversible. We use Metropolis--Hastings
acceptance on cut energy to enforce detailed balance, but the
high-dimensional continuous proposal can have low acceptance rates if
poorly tuned. The mixing-with-ReCom argument guarantees ergodicity
but not necessarily fast mixing.

\textbf{(L5) Integer QP is computationally expensive.} Chain C's MIQP
solves are $\sim 10\times$ slower than Chain B's QP solves, and Chain
B is already $\sim 5\times$ slower than ReCom. For massive precinct
graphs, time limits with warm starts will be necessary.

\subsection{Future directions}\label{sec:disc-future}

\begin{enumerate}
  \item Replica exchange or parallel tempering between Chain B/C and
    Chain V/A to combine compactness with coverage.
  \item Integration with the cycle walk \cite{deford2025cycle} or the
    balanced up-down walk \cite{akitaya2026bud} as alternatives to
    Tempered ReCom in the outer step.
  \item Theoretical mixing-time analysis on simplified grid graphs,
    using the spectral filter form~\eqref{eq:freq-response} as a
    starting point.
  \item Multi-scale diffusion kernels (heat kernel, commute-time
    kernel) and their effect on cut quality vs.\ compactness.
\end{enumerate}

\note{Add a Conclusion section here ($\sim$3 paragraphs) before the
appendix:
(1) one-paragraph summary of the contribution (kernel-aware
local refinement; continuous + integer; 23.9\% cut reduction on Iowa;
narrow-throat diagnostic);
(2) one paragraph on what the framework unlocks for the field (kernel
choice as a first-class modeling decision; the QP/IQP duality as a
general pattern for redistricting MCMC);
(3) one closing paragraph pointing to the most promising follow-up
(integration with reversible / cycle / BUD chains; theoretical mixing
on simplified grids).}

%════════════════════════════════════════════════════════════════════════════════
\appendix
\section{Algorithm details and hyperparameters}\label{app:hyperparams}
%════════════════════════════════════════════════════════════════════════════════

\subsection{Common setup}\label{app:common}

\begin{itemize}
  \item Dual graph $G = (V, E)$ with vertex populations $p_v$ and
    edge attributes (\texttt{shared\_perim} for Iowa).
  \item $k$ districts, $\varepsilon$-population balance per
    Equation~\eqref{eq:pop-balance}.
  \item Initial partition from \texttt{gerrychain.tree.recursive\_tree\_part}.
  \item Step count: $5 \times 5$: $N = 1\,000$, burn-in $B = 100$;
    Iowa: $N = 500$, $B = 100$.
  \item Random seed: 42 (Python \texttt{random} and \texttt{numpy}).
\end{itemize}

\subsection{Chain-specific parameters}\label{app:chains}

\begin{table}[H]
\centering
\small
\begin{tabular}{lll}
\toprule
Chain & Parameter & Value \\
\midrule
A & $\lambda$ (target temperature) & 0.10 \\
B & $p_R$ (ReCom mix probability) & 0.5 \\
B & $\alpha$ (Laplacian weight) & 10.0 \\
B & $\beta$ (fidelity weight) & 1.0 \\
B & Threshold scheme & Cheeger sweep, balance fallback \\
B & Laplacian normalization & False (unnormalized $L$) \\
C & $n_{\mathrm{iqp}}$ (IQP steps per outer) & 10 \\
C & $\alpha$ & 10.0 \\
C & $\beta$ & 1.0 \\
C & MIQP time limit & 5.0 s \\
D & $p_R$ & 0.5 \\
D & Weight randomization & $\mathcal{U}[1, 2]$ per edge \\
D & Repair & False (faithful Davies et al.) \\
\bottomrule
\end{tabular}
\caption{Locked-in hyperparameters for the five
  chains.}\label{tab:hyperparams}
\end{table}

\subsection{Diversity controls}\label{app:diversity}

\begin{itemize}
  \item \textbf{Weight jitter}: multiply each edge weight by
    $\mathcal{U}[1, 2]$ before solving (SpecReCom-style). Tested on B
    and C; gave $+27\%$ coverage gain on C with negligible cut quality
    cost. Always on for D (faithful Davies et al.).
  \item \textbf{Tabu}: forbid the last $K$ pairs chosen by the local
    kernel (cleared whenever ReCom moves). $K = 3$ in our tests;
    modest cut-tightening, no coverage gain.
\end{itemize}

\subsection{Compute environment}\label{app:compute}

\begin{itemize}
  \item macOS (Darwin 23.6.0), Python 3.12.
  \item Gurobi 11 (academic licence) for QP and MIQP.
  \item NetworkX 3.x for graph operations.
  \item Single-threaded; no GPU. All experiments run on a 2024
    MacBook.
\end{itemize}

%════════════════════════════════════════════════════════════════════════════════
\section{Full parameter sweep tables}\label{app:sweeps}
%════════════════════════════════════════════════════════════════════════════════

\needdata{Full sweep tables for $n_{\mathrm{iqp}} \in \{3, 5, 10\}$,
$\alpha \in \{10, 50\}$, $\beta \in \{0, 1\}$, kernel $\in \{k_{\mathrm{unif}},
k_{\mathrm{perim}}, k_{\mathrm{dens}}\}$.}

%════════════════════════════════════════════════════════════════════════════════
\section{Reproducibility checklist}\label{app:reproducibility}
%════════════════════════════════════════════════════════════════════════════════

\begin{itemize}
  \item Code: \texttt{diffusion\_model/recom/} in the project
    repository.
  \item Modules: \texttt{problems.py} (graph setup),
    \texttt{chains.py} (chain runners), \texttt{qp\_model.py}
    (QP/IQP/SpecReCom step factories), \texttt{diagnostics.py}
    (KL/TV diagnostics).
  \item Orchestrator: \texttt{iowa\_target\_matched.py}.
  \item Outputs: \texttt{iowa\_hybrid\_test\_out/target\_matched\_*.npz}.
  \item Fibers: \texttt{fiber\_5x5/} contains the $5 \times 5$ and
    $4 \times 4$ enumeration scripts and pre-computed
    adjacency files.
  \item Random seed: 42. Gurobi default seed (deterministic given
    same input).
\end{itemize}

\bibliographystyle{plain}
\bibliography{ref}

\end{document}
